% !PRINTED 520
\weberpage{520}
%<<CONT
\weq{20}{\begin{aligned}
M &= (-1)^{\frac{m}{8}} \frac{1+i}{\sqrt2} \left(1 - i\tfrac12\sqrt{m}\right)^3, & A &\equiv 1\ (\mathrm{mod}\ 8), \\
M &= (-1)^{\frac{m}{8}} \frac{-1+i}{\sqrt2} \left(1 - i\tfrac12\sqrt{m}\right)^3 = M', & A &\equiv 3\ (\mathrm{mod}\ 8), \\
M &= (-1)^{\frac{m}{8}} \frac{-1-i}{\sqrt2} \left(1 - i\tfrac12\sqrt{m}\right)^3 = M'', & A &\equiv 5\ (\mathrm{mod}\ 8), \\
M &= (-1)^{\frac{m}{8}} \frac{1-i}{\sqrt2} \left(1 - i\tfrac12\sqrt{m}\right)^3 = M''', & A &\equiv 7\ (\mathrm{mod}\ 8).
\end{aligned}}
The four values of $M$ arise from the first through the interchanges:
\weq{21}{\begin{aligned}
M, &\quad \sqrt2, \quad i, \quad \sqrt{m}, \\
M', &\quad -\sqrt2, \quad -i, \quad -\sqrt{m}, \\
M'', &\quad -\sqrt2, \quad i, \quad \sqrt{m}, \\
M''', &\quad \sqrt2, \quad -i, \quad -\sqrt{m}.
\end{aligned}}
If now $\Phi(M,u) = 0$ is, in the same sense as above, the equation holding between $u$ and $M$, then the four functions
\[
\Phi(M,u), \quad \Phi(-M,u), \quad \Phi(iM,u), \quad \Phi(-iM,u)
\]
are without a common divisor, and if $\Phi(M,u)$ is contained in $H(u)$, then the three other functions $\Phi(-M,u)$, $\Phi(iM,u)$, $\Phi(-iM,u)$ are contained in $H(u)$ as well, since the signs of the irrationalities (20) can be changed arbitrarily. It is therefore
\[
H(u) = \Phi(M,u)\,\Phi(-M,u)\,\Phi(iM,u)\,\Phi(-iM,u),
\]
and thus each of the four cases
\[
A \equiv 1, \quad A \equiv 3, \quad A \equiv 5, \quad A \equiv 7\ (\mathrm{mod}\ 8)
\]
occurs in equally many classes of forms of the discriminant $-4m$.

An exception occurs only when $m$ is a square or the double of a square, because in the first case $\sqrt m$ is rational and therefore only the interchange $(M, M'')$ is admitted, while in the second case $\sqrt2$ and $\sqrt m$ change their sign simultaneously, so that only the interchange $(M, M')$ is admissible.

But these interchanges already suffice to split the partial equation further by adjunction of $\sqrt2$.
